%========================================================================================
\section{Conclusion}
\label{sec:conclusion}
%========================================================================================

\subsection{Bilan}

Le back-end de TravelHub répond à tout l'énoncé. Le modèle relationnel a des relations 1-N et une relation N-N avec ses propres attributs, et PostgreSQL vérifie les règles importantes. Les mêmes fonctions sont disponibles dans une API REST documentée avec OpenAPI et dans une API GraphQL. Les deux partagent la logique métier, les permissions et la validation. Les deux services externes sont appelés avec un délai, un nouvel essai, un cache et une dégradation partielle. Pour la sécurité, nous avons suivi l'OWASP API Security Top 10, et 724 tests automatisés couvrent 98\,\% du code.

Quelques choix nous ont beaucoup aidés. Comme toutes les requêtes partent de \code{visible_to()}, le contrôle d'accès au niveau objet est réglé une seule fois, et pas vue par vue. Valider à la fois dans l'application et dans la base nous donne des erreurs claires, tout en gardant les données cohérentes quand deux requêtes arrivent en même temps. Enfin, grâce aux services métier partagés, ajouter l'API GraphQL n'a demandé qu'une nouvelle couche de présentation.

\subsection{Limites}

Notre travail a quand même des limites. En production, nous n'avons pas configuré de moteur de cache. Django utilise donc un cache en mémoire propre à chaque processus. Les réponses externes et les compteurs de limites de requêtes ne sont pas partagés entre les processus Gunicorn. Avec deux processus, la limite réelle est donc deux fois plus haute.

Les prévisions météo ne vont pas au-delà de 16 jours, et Countries donne les noms de pays seulement en anglais. Les risques résiduels de la section~\ref{sec:securite} sont toujours là (on peut deviner des courriels, la largeur des requêtes GraphQL n'est pas limitée, la documentation OpenAPI est publique). Enfin, la pagination GraphQL par décalage devient lente sur de très grandes listes. Une pagination par curseur serait plus solide.

\subsection{Pistes pour le TP2}

Le front-end pourra utiliser directement ce que nous avons fait. GraphQL servira pour les écrans les plus chargés, comme la page d'un voyage avec la météo et le pays, en une seule requête. Il faudra seulement ajouter son origine à \code{CORS_ALLOWED_ORIGINS}. Avec les codes d'erreur stables et le champ \code{details}, on pourra afficher chaque message à côté du bon champ de formulaire. L'en-tête \code{X-Request-ID} permettra aussi de retrouver dans les journaux une erreur vue par l'utilisateur. Côté serveur, un cache partagé comme Redis réglerait la principale limite que nous avons relevée.
